% ============================================================================
%  Chapter 12 -- Verm1, Part 2: Steering, Food and a Score       (Lecture 12)
% ============================================================================
\chapter{Verm1, Part 2: Steering, Food and a Score}
\label{ch:verm1p2}

\epigraphline{A game is a series of interesting choices.}{Sid Meier}

\begin{chapgoals}
By the end of this chapter you will have
\begin{itemize}
\item a worm's head that you steer with the joystick, in four directions;
\item a rule that forbids the U-turn, written as one XOR and one comparison;
\item movement taken from two small tables instead of four branches;
\item food that appears at a random cell, is eaten, and appears again;
\item a score, shown as marks in the title line, and a sound when the worm eats.
\end{itemize}
\end{chapgoals}
% Generated by tools/textbook/export_listings.js from tools/decks/mastery.js and tools/concepts/registry.json: do not edit.
\begin{chapstatus}{Learning status of this chapter}
\masterychip{now}{UNDERSTAND NOW}\enspace bit masks; codes; arrays in memory; pull-up inputs; Boolean algebra; fixed point; multiplexers; signed or unsigned instructions; comparators and parity; arguments in a0 to a7.\par\smallskip
\masterychip{use}{USE NOW}\enspace the calling convention (Chapter~\ref{ch:processor}) and, used without being studied in this course, the PRG32 tools, the PRG32 runtime calls, the PRG32-QT debugger.\par\smallskip
\masterychip{later}{UNDERSTAND LATER}\enspace nothing.
\end{chapstatus}
\begin{chapstatus}{Roadmap of this chapter}
By the end of the chapter you can:
\begin{enumerate}[nosep,leftmargin=2em]
\item read the buttons with masks and refuse a U-turn with one XOR
\item compare positions with XOR and OR to detect food
\item keep and show a score
\end{enumerate}
It builds on: bit masks (Chapter~\ref{ch:boole1}); codes (Chapter~\ref{ch:combinational}); arrays in memory (Chapter~\ref{ch:information}); pull-up inputs (Chapter~\ref{ch:electronics}); Boolean algebra (Chapter~\ref{ch:boole1}). Next: Chapter~\ref{ch:sequential}.
\par\smallskip\textbf{Looking back, looking ahead.} In Lecture 07 a function that calls another saved ra in a frame of 16 bytes. Today the functions of Verm1 call each other more deeply: each saves ra once, on entry, and restores it before ret. Lecture 18 states the rules: which registers a call may change and which it must preserve. \emph{Recall:} A function calls two other functions. How many times must it save ra, and where?
\end{chapstatus}


In this second tutorial the player enters the game. Every idea used has
already been seen as a circuit in Chapters \ref{ch:boole1}
to~\ref{ch:combinational}; here the same ideas are instructions
(Table~\ref{tab:v2:from}). The reference file is \fname{verm1\_part2.S}
(Appendix~\ref{app:verm1}).

\begin{table}[htb]
\centering\small
\caption{What Part 2 takes from the chapters on logic.}\label{tab:v2:from}
\begin{tabularx}{\textwidth}{@{}lLL@{}}
\toprule
\textbf{From} & \textbf{Idea} & \textbf{Where it appears}\\
\midrule
Chapter \ref{ch:boole1} & AND as a mask: keep one bit, clear the others & one button out of the word of buttons\\
Chapter \ref{ch:boole1} & truth tables: a code for each case & the four directions in two bits\\
Chapter \ref{ch:boole2} & XOR is 1 where two bits differ & the U-turn rule; the test for equality\\
Chapter \ref{ch:electronics} & a button is a switch to ground with a pull-up & why ``pressed'' arrives as 1 only after the runtime inverts it\\
Chapter \ref{ch:combinational} & comparator, multiplexer, ROM as a table & head on the food; \code{dx} and \code{dy}\\
\bottomrule
\end{tabularx}
\end{table}

\section{Start from a Part 1 that runs}

Open your \fname{verm1.S} of Part~1, build it and run it with the two
commands of Listing~\ref{lst:v1:build}: the head must cross the screen. If it
does not, copy the reference file \fname{verm1\_part1.S} and start from that.
The file is changed in three steps and is run after the second and after the
third. Listings marked \textbf{changed} replace code that the file already
contains; listings marked \textbf{new} add code. Keep a copy of each finished
part, and keep the debugger of \prgqt open: step~4 uses it to watch a decision
being taken.

% ----------------------------------------------------------------------------
\section{Step 1: names for bits}

\subsection{Constants: buttons and directions}

Add the lines from \code{RED} and from \code{BTN\_LEFT} to \code{DIR\_UP} to
the constants of Part~1 (Listing~\ref{lst:v2:const}). There is a new colour,
red, for the food. Each button is one bit of a word, and a mask selects it. A
direction is a number from 0 to 3: two bits.

\asmfile[caption={Step 1 \textbf{(changed)}: constants.},label=lst:v2:const]{code/l12/c04.S}

\subsection{A code chosen with care}

The codes of the four directions are not arbitrary
(Table~\ref{tab:v2:dirs}). Bit~0 says ``vertical'' and bit~1 says
``backwards'', so that a direction and its opposite differ in bit~1 and only
there: the XOR of a code with its opposite is always $10_2=2$.

\begin{table}[htb]
\centering\small
\caption{The four directions. On the screen $y$ grows downwards: DOWN is
$+1$, UP is $-1$.}\label{tab:v2:dirs}
\begin{tabular}{@{}lcccrrl@{}}
\toprule
\textbf{Direction} & \textbf{Code} & \textbf{Bit 1} & \textbf{Bit 0} & \textbf{dx} & \textbf{dy} & \textbf{Opposite}\\
\midrule
RIGHT & 0 & 0 & 0 & $+1$ & 0 & LEFT = 2\\
DOWN & 1 & 0 & 1 & 0 & $+1$ & UP = 3\\
LEFT & 2 & 1 & 0 & $-1$ & 0 & RIGHT = 0\\
UP & 3 & 1 & 1 & 0 & $-1$ & DOWN = 1\\
\bottomrule
\end{tabular}
\end{table}

\begin{prgconcept}
A good encoding turns a rule into one gate. It is the lesson of
Chapter~\ref{ch:boole2}, and it is worth more than any trick of programming.
\end{prgconcept}

\subsection{Data: two tables and five new variables}

Replace the data of Part~1 with Listing~\ref{lst:v2:data}. The tables
\code{dx} and \code{dy} have four words each, one per direction; they are
constants, and therefore live in \code{.rodata}. The variable \code{dir} is
the direction of the last step, \code{want} the one asked for by the player;
\code{score}, \code{food\_x} and \code{food\_y} are used in step~3.

\asmfile[caption={Step 1 \textbf{(changed)}: data.},label=lst:v2:data]{code/l12/c05.S}

% ----------------------------------------------------------------------------
\section{Step 2: steering}

The plan is: read the buttons, refuse the U-turn, look the movement up in a
table.

\subsection{From finger to bit}

The pin of a button reads 0 when the button is pressed: it is active-low
(Chapter~\ref{ch:electronics}). The function \code{prg32\_input\_read} inverts
and packs: one word, one bit per button, with $1=\text{pressed}$. \key{UP}
and \key{RIGHT} together give $\hex{04}+\hex{02}=\hex{06}$, and
$\reg{a0}\cdot\hex{04}$ is zero if and only if \key{UP} is not pressed. The
cartridge never sees the electrical level: the runtime gives the same word on
the board, on QEMU and on \prgqt.

\subsection{\texttt{update}: from buttons to a wish}

\asmfile[caption={Step 2 \textbf{(changed)}: the beginning of \code{verm1\_update}.},label=lst:v2:wish]{code/l12/c06.S}

Listing~\ref{lst:v2:wish} replaces the beginning of \code{verm1\_update},
which now calls a function and therefore needs a stack frame. \ins{andi}
keeps one bit, and zero means ``not pressed''. Each pressed button overwrites
the wish in \reg{t2}; with two buttons pressed the last test wins, so that
\key{UP} beats \key{LEFT}, which beats \key{DOWN}, which beats \key{RIGHT}.
Then comes the rule: the XOR of the wish and of the current direction is 2
exactly when they are opposite, and in that case the wish is not stored and
the old one survives.

The wish is stored in \code{want} and becomes the direction only at the next
step. In this way two quick presses between one step and the next cannot fold
the worm back onto itself.

\begin{table}[htb]
\centering\small
\caption{The U-turn rule as a truth table: six of the sixteen rows.}\label{tab:v2:uturn}
\begin{tabular}{@{}llccl@{}}
\toprule
\textbf{dir} & \textbf{want} & \textbf{dir $\oplus$ want} & \textbf{$=2$?} & \textbf{Decision}\\
\midrule
0 RIGHT & 0 RIGHT & \texttt{00} & no & go on\\
0 RIGHT & 1 DOWN & \texttt{01} & no & turn\\
0 RIGHT & 2 LEFT & \texttt{10} & yes & refuse: U-turn\\
0 RIGHT & 3 UP & \texttt{11} & no & turn\\
1 DOWN & 3 UP & \texttt{10} & yes & refuse: U-turn\\
3 UP & 2 LEFT & \texttt{01} & no & turn\\
\bottomrule
\end{tabular}
\end{table}

The rule is one XOR and one comparison with the constant 2
(Table~\ref{tab:v2:uturn}). Complete the sixteen rows as an exercise and
count the refusals: there are four, one per direction. Without the careful
encoding the rule would need four comparisons.

\subsection{\texttt{update}: the timer calls \texttt{step}}

\asmfile[caption={Step 2 \textbf{(changed)}: the rest of \code{verm1\_update}.},label=lst:v2:timer]{code/l12/c07.S}

The fixed-point timer of Part~1 is unchanged (Listing~\ref{lst:v2:timer}).
The lines that moved \code{head\_x} are gone: moving is now a function of its
own, \code{step}. The local labels are 6 and 7 because the numbers 1 to 5 are
used above.

\subsection{\texttt{step}: look the movement up}

\asmfile[caption={Step 2 \textbf{(new)}: \code{step}, first half.},label=lst:v2:step1]{code/l12/c08.S}

The wish becomes the direction, once per step
(Listing~\ref{lst:v2:step1}). Then the direction selects the movement: a word
is 4 bytes, so the offset is the code shifted to the left by 2, and the
address of the table plus the offset is the address of the word to load.
There is no branch: the code selects, as in a multiplexer. A table in
read-only memory addressed by a code is a ROM (Chapter~\ref{ch:storage}). The
alternative is four branches; the table is shorter, and changing the rules
means changing data, not code.

\subsection{\texttt{step}: move and wrap, both ways}

\asmfile[caption={Step 2 \textbf{(new)}: \code{step}, second half.},label=lst:v2:step2]{code/l12/c09.S}

In Part~1 only the right edge existed. Now \code{dx} can be $-1$ and the
column can go below 0: a negative column is a signed number, and \ins{bgez}
and \ins{blt} are the signed comparisons of Chapter~\ref{ch:integers}
(Listing~\ref{lst:v2:step2}). The same is done for the row, with
\code{ROWS}. The registers \reg{t4} and \reg{t6} keep the new cell for step~3.

For the moment the function is closed with the three lines of
Listing~\ref{lst:v2:end}, which restore \reg{ra} and return. Step~3 will put
the test for the food before them.

\asmfile[caption={Step 2: a temporary ending for \code{step}.},label=lst:v2:end]{code/l12/c10.S}

\begin{prgcheck}[Run it: you are driving]
\begin{enumerate}[nosep]
\item Build and run: the head still starts to the right, from column 5 of
  row 12.
\item Hold \key{DOWN}: the head turns at its next step and goes down, wrapping
  from row 23 to row 0.
\item While it goes right, hold \key{LEFT}: nothing happens. The U-turn is
  refused.
\item Hold \key{DOWN} and then \key{LEFT}: two turns, and now the head goes
  left.
\end{enumerate}
\textbf{Expected result.} Holding \key{DOWN} from the start for 64 frames
brings the head to column 5, row 22: ten steps, all downwards. Holding
\key{LEFT} for 64 frames brings it to column 15, row 12, as if nothing had
been pressed. (Both were checked on the headless runner of \prgqt.)
\end{prgcheck}

% ----------------------------------------------------------------------------
\section{Step 3: food, a score, a sound}

\subsection{\texttt{init}: a complete start}

\asmfile[caption={Step 3 \textbf{(changed)}: \code{verm1\_init}.},label=lst:v2:init]{code/l12/c12.S}

Every variable now receives its starting value
(Listing~\ref{lst:v2:init}). \code{init} calls a function, to place the first
food, and so it needs a stack frame.

\subsection{\texttt{place\_food}: a random cell}

\asmfile[caption={Step 3 \textbf{(new)}: \code{place\_food}.},label=lst:v2:food]{code/l12/c13.S}

The function \code{prg32\_random\_number(min, max)} returns a number between
its two arguments, both included; the arguments go in \reg{a0} and \reg{a1}
and the result comes back in \reg{a0} (Listing~\ref{lst:v2:food}). The random
numbers come from the runtime, and on the board from the hardware generator of
the ESP32-C6: two runs of the game are different. The food may fall on the
head; Part~3 fixes it.

\subsection{\texttt{step}: is the head on the food?}

\asmfile[caption={Step 3 \textbf{(changed)}: the end of \code{step}.},label=lst:v2:eat]{code/l12/c14.S}

Replace the temporary ending with Listing~\ref{lst:v2:eat}. A comparator for
equality is a row of XOR gates followed by an OR
(Chapter~\ref{ch:combinational}); here it is two \ins{xor} and one \ins{or}:
each \ins{xor} is 0 when two numbers are equal, and the \ins{or} of the two
results is 0 only if both pairs are equal. Eating means one more point, new
food and a sound. The registers \reg{t4} and \reg{t6} still hold the new
column and row; after the calls they are no longer needed, which is why it is
safe to keep them in temporaries.

\subsection{\texttt{eat\_sound}: one note}

\asmfile[caption={Step 3 \textbf{(new)}: \code{eat\_sound}.},label=lst:v2:sound]{code/l12/c15.S}

The signature is
\begin{center}\code{prg32\_audio\_note(channel, instrument, note, volume, milliseconds)}\end{center}
with five arguments, \reg{a0} to \reg{a4}, then the call. Note 72 is C5 in the
MIDI numbering. On \prgqt the note is played by the computer, on the board by
the amplifier and its speaker. Sound comes last: the game worked without it.

\subsection{\texttt{draw}: the score as marks}

\asmfile[caption={Step 3 \textbf{(changed)}: the beginning of \code{verm1\_draw}.},label=lst:v2:marks]{code/l12/c16.S}

Listing~\ref{lst:v2:marks} draws one small rectangle per point, 4 pixels
apart, in a loop that ends when the counter reaches the score. The counter is
kept in \reg{s0}, because it must survive the calls: a \reg{t} register may be
destroyed by any call, whereas an \reg{s} register is preserved, provided
that the function which uses it saves it and restores it. So \reg{s0} is
saved in the frame, with \reg{ra}. The convention, met in
Chapter~\ref{ch:firstcontact}, is studied in Chapter~\ref{ch:processor}. Digits arrive in Part~4, when division is
available.

\subsection{\texttt{draw}: food and head}

\asmfile[caption={Step 3 \textbf{(changed)}: the rest of \code{verm1\_draw}.},label=lst:v2:draw]{code/l12/c17.S}

The same three lines appear twice, for column, row and colour
(Listing~\ref{lst:v2:draw}); the conversion to pixels is left to a function.
The head is drawn last, so that it covers the food.

\asmfile[caption={Step 3 \textbf{(new)}: \code{draw\_cell}.},label=lst:v2:cell]{code/l12/c18.S}

\code{draw\_cell} is the arithmetic of Part~1, written once
(Listing~\ref{lst:v2:cell}). It is a function with a contract:
$\reg{a0}=\text{column}$, $\reg{a1}=\text{row}$, $\reg{a2}=\text{colour}$.
The colour arrives in \reg{a2} and must leave in \reg{a4}: it is moved first,
because \reg{a2} is about to be overwritten. Writing the contract in a
comment is part of the job.

\begin{prgcheck}[Run it: Verm1, Part 2]
A red cell is somewhere on the screen. Steer the head onto it: a note sounds,
a mark appears next to \texttt{VERM1}, the food moves
(Figure~\ref{fig:v2:shot}).
\end{prgcheck}

\begin{figure}[htb]
\centering
\includegraphics[width=0.55\textwidth]{verm1_p2.png}
\caption{The head, the food, and two marks after the title: two points. The
screenshot is produced by the tests, in which a scripted player steers the
head to the food twice.}
\label{fig:v2:shot}
\end{figure}

% ----------------------------------------------------------------------------
\section{Step 4: watch a decision}

In this session \prgqt is used as a logic probe.

\begin{prgtry}[Catch the U-turn rule at work]
\begin{enumerate}[nosep]
\item Run the cartridge in debug mode and press \key{Update}: the listing shows
  \code{verm1\_update}.
\item Look at the telemetry while you hold a button: the input mask shows
  \hex{04} for \key{UP}, \hex{01} for \key{LEFT}.
\item Find the instruction \code{xor t3, t2, t1} and click in the gutter: a
  breakpoint.
\item Hold \key{LEFT} while the head goes right, and press \key{Resume}: the
  program stops on the \ins{xor}.
\item Read the registers: \reg{t1} is 0, the direction; \reg{t2} is 2, the
  wish.
\item Press \key{Step}: \reg{t3} becomes 2. Step twice more: the branch is
  taken and the \ins{sw} is skipped.
\end{enumerate}
Lower the speed of the debugger to 0.25x if you need time to press.
\end{prgtry}

Table~\ref{tab:v2:regs} lists what is read in the registers, always with the
head going right. Its last row shows the priority of the tests: \key{UP} is
tested last and wins. \reg{a0} is overwritten later in the frame, and must be
read at the first instruction after the call.

\begin{table}[htb]
\centering\small
\caption{What the registers show at the breakpoint.}\label{tab:v2:regs}
\begin{tabular}{@{}lccccl@{}}
\toprule
\textbf{You hold} & \textbf{\reg{a0}} & \textbf{\reg{t1} (dir)} & \textbf{\reg{t2} (want)} & \textbf{\reg{t3} = \reg{t2}$\oplus$\reg{t1}} & \textbf{Stored?}\\
\midrule
nothing & \hex{00} & 0 & 0 & 0 & yes: 0\\
\key{DOWN} & \hex{08} & 0 & 1 & 1 & yes: 1\\
\key{LEFT} & \hex{01} & 0 & 2 & 2 & no: U-turn\\
\key{UP} & \hex{04} & 0 & 3 & 3 & yes: 3\\
\key{UP} + \key{RIGHT} & \hex{06} & 0 & 3 & 3 & yes: 3\\
\bottomrule
\end{tabular}
\end{table}

\subsection{Read the variables in memory}

After \code{la t0, dir} the register \reg{t0} holds the address of
\code{dir}: type it in the \emph{Memory} field. The variables follow each
other, 4 bytes each, in the order of the file: from the address of \code{dir},
at $+0$ there is \code{dir}, at $+4$ \code{want}, at $+8$ \code{step\_q8}, at
$+12$ \code{score}, at $+16$ \code{food\_x} and at $+20$ \code{food\_y};
\code{head\_x} and \code{head\_y} are the two words before \code{dir}. Eat
once, pause, and read \code{score}: \texttt{01 00 00 00}, lowest byte first.
Then read \code{food\_x} and \code{food\_y} and check them against the screen:
they are the column and the row of the red cell.

% ----------------------------------------------------------------------------
\section{Check your understanding}

\begin{quiz}{The player holds \key{DOWN} and \key{RIGHT} together: \code{prg32\_input\_read} returns}
\item \hex{02}
\item \hex{08}
\item \hex{0a}
\item \hex{10}
\end{quiz}

\begin{quiz}{The worm goes UP (3). Which wish is refused?}
\item RIGHT (0)
\item DOWN (1)
\item LEFT (2)
\item UP (3)
\end{quiz}

\begin{quiz}{Why \code{slli t1, t1, 2} before adding to the address of \code{dx}?}
\item to make the code positive
\item each entry of the table is 4 bytes
\item to clear bit 1
\item tables start at address 4
\end{quiz}

\begin{quiz}{\code{xor t1, t1, t4} gives zero exactly when}
\item \reg{t1} is zero
\item \reg{t4} is zero
\item \reg{t1} and \reg{t4} are equal
\item \reg{t1} and \reg{t4} are opposite
\end{quiz}

\section{Make it yours}

\begin{exercise}[on \prg]
\begin{enumerate}[nosep]
\item Make button \key{A} a brake: while it is held, the timer does not advance.
\item Give two points for a piece of food eaten in the top row.
\item Draw the head in a different colour for each direction, using a third
  table of four colours.
\item Change the encoding so that RIGHT = 0, LEFT = 1, DOWN = 2, UP = 3: what
  does the U-turn rule become?
\end{enumerate}
\expected{Each change is a few lines. Keep a copy of the working Part~2
before experimenting: Part~3 starts from it. (Hints: for the brake,
\ins{andi} with \code{BTN\_A}, \hex{10}, at the start of the timer and a
branch to the exit. The table of colours is four \code{.word} in
\code{.rodata}, indexed like \code{dx}. With the new encoding the opposites
differ in bit~0, so the test becomes $\oplus = 1$.)}
\end{exercise}

\section{Where we are, and what comes next}

The player is in the game: a wish, a rule, a movement, a reward. Bits and
masks, XOR as a comparator, a table as a selector: logic has been turned into
instructions. But the worm has no body, nothing can go wrong, and the game
never ends. Chapters \ref{ch:sequential} to~\ref{ch:storage} give the tools:
state machines, registers, memory as an array, and what survives when the
power is switched off. Part~3 adds a body that grows, collisions, a title
screen and a game over.
